\chapter{What this toolkit reaches}

This chapter assesses which problem an arbitrary-precision apparatus could be pointed at, the
reasoning for it, the trap beside it, and the measured bound that constrains it. The assessment is a
recommendation and no work on it is begun.

\section{The obvious answer, and why it is the trap}

An apparatus whose edge is arbitrary precision reaches for Riemann first, and Riemann is the wrong
choice for a reason that has nothing to do with how good the arithmetic is.

The apparatus matches the \emph{subject} and not the \emph{obstruction}. A zeta function and
unbounded precision look like a fit. But zeros on the critical line are computed far past anything a
desktop reaches. The first $10^{13}$ zeros are located on the line~\cite{src:Gourdon-2004}, zeros
near the $10^{22}$nd are computed~\cite{src:Odlyzko-2001}, and the hypothesis is verified rigorously
up to height $3\cdot10^{12}$~\cite{src:Platt-and-Trudgian-2021}. No finite count of zeros bears on
the proof at all. Adding zeros reproduces existing work at a worse scale. The conjecture is not
precision-bound, and an apparatus whose advantage is precision buys nothing against it.

That is the superficial-match failure in its exact form, and it is worth naming because it is the
shape of mistake this apparatus is most exposed to: the subject looks like the instrument.
Yang--Mills wants constructive quantum field theory and lattice gauge machinery. Hodge wants
algebraic geometry with no numerical handle this toolkit has. P versus NP is not a numerical
question.

\section{The recommendation}

Birch and Swinnerton-Dyer, because its obstruction matches the instrument, where the other
Millennium problems match only in subject.

It is the only one of them found by computation: Birch and Swinnerton-Dyer fitted curves numerically
on EDSAC~\cite{src:Birch-and-Swinnerton-Dyer-1965}. It is natively experimental, and the refined
conjecture is an exact numerical identity~\cite{src:Wiles}. For an elliptic curve $E$ of rank $r$,
writing $\mathrm{Sha}$ for the order of the Tate--Shafarevich group, $\Omega$ for the real period,
$\mathrm{Reg}$ for the regulator and $c_p$ for the Tamagawa numbers,
\[
  \frac{L^{(r)}(E,1)}{r!} \;=\; \frac{\mathrm{Sha}\cdot\Omega\cdot\mathrm{Reg}\cdot\prod_p c_p}
                                    {|E_{\mathrm{tors}}|^2}.
\]
Rearranged, the analytic order of $\mathrm{Sha}$ is a quotient of computable real numbers, and the
conjecture asserts that this quotient is an integer, usually a perfect square. The verification is
therefore: compute a real number to $N$ digits and ask whether it lands on an integer. That is a
precision-bound question, and it is the check the arithmetic already performs.

\section{The first measurement, and the positive control it requires}

The positive control comes first and is not optional. Curves are taken from published
tables~\cite{src:Cremona-1997} where the rank and the order of $\mathrm{Sha}$ are already known, at
rank zero and one, small conductor, with $\mathrm{Sha}$ trivial or known to be 4 or 9. The quotient
is computed and must land on the known integer. An instrument that cannot reproduce a known answer
says nothing when pointed at an unknown one.

\textbf{The quotient is only as good as its weakest input.} The real period and the regulator come
from their own algorithms carrying their own error. The report is therefore the precision of the
quotient and not the precision of the arithmetic. An arithmetic with no floor under it does not give
its inputs one.

\section{The null, drawn and not derived}

\textbf{The derived null does not hold.} It assumes that a real number carrying no integrality
structure has its distance to the nearest integer distributed uniformly on $[0,\tfrac12]$, that
landing within $10^{-N}$ of an integer by chance therefore has probability about $2\times10^{-N}$,
and that agreement to fifty digits is a bound of $10^{-50}$.

\textbf{What refutes it.} It is derived and not drawn. A threshold reasoned out from an assumed
distribution omits the variance sources not enumerated, and omitted variance only ever pushes one
way.

One narrower objection to it is itself wrong. It holds that the integer factors in the quotient, the
torsion order squared, the Tamagawa product and $r!$, might concentrate the quotient near integers
and inflate the coincidence rate. They do not: scaling a structureless real by a known rational does
not concentrate it near integers. The derived null fails for the better reason and not for that one.

\textbf{The drawn null} measures the null under the same conditions as the effect, through the same
pipeline and with the same inputs and the same precision losses, in two ways. Forming the quotient
with the rank taken one too high or one too low gives a condition under which the identity should
not hold, and the distribution of distance-to-nearest-integer across many curves is then the null
itself. Computing the quotient across many curves whose $\mathrm{Sha}$ is known gives the spread of
the residual about the true integer. That spread is the instrument's own floor, and the period and
the regulator dominate it well before the arithmetic does. The second is the positive control and
the floor measurement in one pass.

The consequence changes what may be claimed. The bound is not $10^{-N}$ for $N$ digits of
arithmetic. It is however close the pipeline gets on a curve whose answer is already known. If the
regulator runs out at $10^{-12}$, then $10^{-12}$ is the claim however many digits the multiply
carries.

\section{The measured bound that constrains this}

The cost of depth in the integer relation search of the same
arithmetic~\cite{src:Quigg-instruments} is measured, and the measurement constrains the plan.

Cost goes as precision to the power 2.64. The prediction is 2.585, from a Karatsuba host multiply
near width$^{1.585}$ against a count linear in the width, and the reduction step does not account
for the difference.

The device overtakes the host at 1{,}024 limbs, 32{,}768 bits. The mean operand runs at 0.642 of the
precision, and the mean multiply reaches that crossover near 51{,}000 binary places. One cell at
51{,}036 places costs 26.2 hours.

None of that binds the recommended measurement, which wants tens to hundreds of digits. It binds two
other things: relation detection at tens of thousands of digits is out of reach, and so is any plan
that relies on the device to rescue it.

\section{What is not claimed}

This is verification and not proof, and that distinction governs everything claimed in this
chapter.

A measurement of this kind produces evidence for particular instances of the conjecture, in the way
the original EDSAC work did. A verified instance at unusual rank or unusual precision is a real
contribution and it is not a solution.

The recommendation carries plain caveats. No expertise in elliptic-curve $L$-function computation
stands behind it. Computing $L^{(r)}(E,1)$ requires modular symbols~\cite{src:Cremona-1997} or
Dokchitser's algorithm~\cite{src:Dokchitser-2004}, established systems already implement those, and
rebuilding them would be wasted effort. The realistic shape is established algorithms for the
$L$-function with this arithmetic supplying precision underneath, and the first task is finding out
whether that seam works at all.
